\chapter*{Preface}
\usetikzlibrary{fit}
\addcontentsline{toc}{chapter}{Preface}

This book combines three sets of lecture notes by Richard Melrose:
\emph{Introduction to Microlocal Analysis} (MIT, version 0.5B, 2003), a revised version of
the notes from several courses he taught at MIT; \emph{Lectures on
Pseudodifferential Operators} (MIT course 18.157, Fall 2005, revised 2006); and
\emph{From Microlocal to Global Analysis} (MIT, version 0.7I, 2008), a later expansion of
the first.

Microlocal analysis is a geometric theory of distributions, or equivalently a theory of
geometric distributions. Rather than general distributions, which are like general
continuous functions only worse, it studies the distributions that actually arise from
differential and integral equations. Conormal distributions and, more generally,
Lagrangian distributions are smooth away from a well-defined set and have a
precisely described singularity there. Their singularities are the main subject.

\subsection*{Prerequisites and suggested reading}
The reader should already know some distribution theory and some basic analysis, and
some familiarity with analysis on manifolds will help. If one or two of these are
missing they can be picked up along the way. A reader who has none of them should first
do some preliminary reading. Friedlander's book \cite{Friedlander} is a good
introduction to distributions, and Volume~I of H\"ormander's treatise
\cite{Hormander-I} gives a much fuller treatment that also serves as a general
reference. For analysis on manifolds, most of what we need can be found in Munkres'
book \cite{Munkres} or in Spivak's short book \cite{Spivak}. \Cref{ch:distributions}
reviews tempered distributions and the Fourier transform, so that the book is
reasonably self-contained.

\subsection*{About this edition}
This is a streamlined edition. The material from the three sources has been merged and
reorganized into five parts.
\begin{itemize}[nosep]
\item Part~I builds the calculus: tempered distributions and compactifications,
  pseudodifferential operators on $\R^n$, microlocalization, conormal distributions and
  operators on manifolds, and the semiclassical and adiabatic calculus.
\item Part~II covers the global calculi on $\R^n$: the isotropic calculus, and traces and
  determinants.
\item Part~III covers boundary problems and propagation: boundary problems, the
  Lagrangian distributions, the Hamilton--Jacobi and transport equations, the wave
  equation and Weyl asymptotics.
\item Part~IV is index theory, taken mainly from the 18.157 notes: homotopy invariance
  of the index, Toeplitz operators, K-theory and Bott periodicity, product-type operators,
  the families index theorem (with a second proof by semiclassical and adiabatic limits,
  from the 2008 notes) and eta forms.
\item Part~V covers the algebraic structure of the calculus: Hochschild homology.
\end{itemize}
Notation has been simplified and made uniform throughout. Proofs that were only
sketched, or left as exercises but needed later, have been completed, and errors have
been corrected. The ``addenda'' to the 18.157 lectures have been moved into the main text
or the exercises. Draft material has been removed, including chapters that existed in
the originals only as headings. Each chapter ends with exercises and a short note on its
main departures from the original notes.

\subsection*{How to read this book}
The book is written for readers at the beginning of graduate school. Each chapter opens
by saying what problem it solves and how it connects to what came before. Short
\emph{Discussion} paragraphs explain the idea behind a definition or a proof, and
\emph{Question} and \emph{Answer} pairs take up points that often puzzle a first-time
reader. These can be skipped by readers who already know the material.

\Cref{preface:fig:dependencies} shows how the chapters depend on each other. Part~I is
the foundation for everything else. The other four parts depend on Part~I but are
largely independent of each other. The main exception is Part~IV: the proof of Bott
periodicity in \cref{ch:bott} uses the isotropic calculus of \cref{ch:isotropic}, and
traces and determinants from \cref{ch:determinants} appear occasionally. Within each
of Parts~III and~IV the chapters are best read in order.

For a first one-semester course we suggest the following path through Part~I.
\begin{itemize}[nosep]
\item \Cref{ch:distributions}, possibly quickly if distributions are familiar. The
  radial compactification (\cref{distributions:sec:radial}) is used throughout and
  should not be skipped; \cref{distributions:sec:other} can wait.
\item \Cref{ch:euclidean} up to and including Sobolev spaces
  (\cref{euclidean:sec:sobolev}). The later sections on topologies and holomorphic
  families are needed only for complex powers and determinants.
\item \Cref{ch:microlocal}: wavefront sets, the characteristic variety and microlocal
  elliptic regularity, followed by propagation of singularities
  (\cref{microlocal:sec:propagation}).
\item \Cref{ch:manifolds} up to the Fredholm theory of elliptic operators
  (\cref{manifolds:sec:fredholm}). The facts about conormal distributions that it uses
  are recalled in \cref{manifolds:sec:recall}, so \cref{ch:conormal} can be read
  lightly on a first pass.
\end{itemize}
A second course can then follow either the analytic route (Part~III) or the
topological route (Part~II and Part~IV).

\begin{figure}[htbp]
\centering
\begin{tikzpicture}[>=Latex,
  ch/.style={draw,rounded corners=2pt,align=center,font=\footnotesize,
    text width=2.75cm,minimum height=0.9cm,inner sep=2pt,fill=white},
  part/.style={draw=gray,dashed,rounded corners=4pt,inner sep=5pt},
  plab/.style={font=\footnotesize\bfseries,anchor=south west,inner sep=1pt}]
  % Part I
  \node[ch] (dist) at (0,0) {\labelcref{ch:distributions}. Distributions};
  \node[ch] (eucl) at (0,-1.3) {\labelcref{ch:euclidean}. Operators\\on $\R^n$};
  \node[ch] (micro) at (0,-2.6) {\labelcref{ch:microlocal}. Microlocalization};
  \node[ch] (conormal) at (0,-3.9) {\labelcref{ch:conormal}. Conormal\\distributions};
  \node[ch] (manifolds) at (0,-5.2) {\labelcref{ch:manifolds}. Operators\\on manifolds};
  \node[ch] (semicl) at (0,-6.5) {\labelcref{ch:semiclassical}. Semiclassical\\calculus};
  % Part II
  \node[ch] (iso) at (4.7,0) {\labelcref{ch:isotropic}. Isotropic and\\scattering calculi};
  \node[ch] (det) at (7.9,0) {\labelcref{ch:determinants}. Traces and\\determinants};
  % Part IV
  \node[ch] (index) at (4.7,-1.9) {\labelcref{ch:index}. Index and\\Toeplitz operators};
  \node[ch] (bott) at (7.9,-1.9) {\labelcref{ch:bott}. K-theory and\\Bott periodicity};
  \node[ch] (product) at (11.1,-1.9) {\labelcref{ch:product}. Product-type\\operators};
  \node[ch] (families) at (11.1,-3.2) {\labelcref{ch:families}. Families\\index theorem};
  \node[ch] (semiidx) at (7.9,-3.2) {\labelcref{ch:semiclassical-index}. Semiclassical\\index theorem};
  \node[ch] (eta) at (11.1,-4.5) {\labelcref{ch:eta}. Eta forms};
  % Part III
  \node[ch] (boundary) at (4.7,-6.5) {\labelcref{ch:boundary}. Elliptic\\boundary problems};
  \node[ch] (hj) at (7.9,-6.5) {\labelcref{ch:hj}. Lagrangian\\distributions};
  \node[ch] (wave) at (11.1,-6.5) {\labelcref{ch:wave}. Wave equation,\\Weyl law};
  % Part V
  \node[ch] (hoch) at (4.7,-8.3) {\labelcref{ch:hochschild}. Hochschild\\homology};
  % part boxes
  \node[part,fit=(dist)(semicl),label={[plab]north west:Part I}] (PI) {};
  \node[part,fit=(iso)(det),label={[plab]north west:Part II}] (PII) {};
  \node[part,fit=(index)(product)(eta),label={[plab]north west:Part IV}] (PIV) {};
  \node[part,fit=(boundary)(wave),label={[plab]north west:Part III}] (PIII) {};
  \node[part,fit=(hoch),label={[plab]north west:Part V}] (PV) {};
  % arrows inside Part I
  \foreach \a/\b in {dist/eucl,eucl/micro,micro/conormal,conormal/manifolds,manifolds/semicl}
    \draw[->] (\a)--(\b);
  % Part II, III, IV
  \draw[->] (iso)--(det);
  \draw[->] (hj)--(wave);
  \draw[->] (index)--(bott); \draw[->] (bott)--(product); \draw[->] (product)--(families);
  \draw[->] (families)--(semiidx); \draw[->] (families)--(eta);
  % between parts
  \draw[->,thick,gray] (PI.east|-iso) -- (PII.west|-iso);
  \draw[->,thick,gray] (PI.east|-bott) -- (PIV.west|-bott);
  \draw[->,thick,gray] (PI.east|-boundary) -- (PIII.west|-boundary);
  \draw[->,thick,gray] (PI.south) |- (PV.west);
  \draw[->] (iso.south east) -- (bott.north west);
  \draw[->,dashed] (semicl.east) .. controls (3,-5.6) and (5.5,-3.9) .. (semiidx.south west);
\end{tikzpicture}
\caption{Dependencies between the chapters. Grey arrows mean that a part uses Part~I;
black arrows show the main dependencies between individual chapters. The dashed arrow
indicates that the semiclassical proof of the index theorem in
\cref{ch:semiclassical-index} uses the semiclassical calculus of
\cref{ch:semiclassical}.}
\label{preface:fig:dependencies}
\end{figure}

\subsection*{Acknowledgements}
Melrose thanks many people for their help with the notes, including Benoit Charbonneau,
Sine Rikke Jensen, Edith Mooers, Mark Joshi, Raul Tataru and Jonathan Kaplan, and, for
the 18.157 lectures, Pierre Albin, Zuoqin Wang, Ben Mares and Austin Ford.
